% einkommentieren für Buchbindung mit zweiseitigem Druck:
% \documentclass[11pt, twoside,openright] {scrreprt}

% Einseitige Version für normalen Ausdruck und die digitale Version
\documentclass[11pt] {scrreprt}

%Papiermaße
\setlength\paperheight {297mm}
\setlength\paperwidth  {210mm}

% Choose your language
%\usepackage{ngerman}
%\usepackage[ngerman]{babel}
\usepackage[english]{babel}

\usepackage{a4,german}
\usepackage[utf8]{inputenc}
\usepackage[T1]{fontenc}

\usepackage{amsfonts}
\usepackage{amsmath}
\usepackage{amssymb}
\usepackage{fancyhdr}
\usepackage{graphicx}
\usepackage{wrapfig}
\usepackage{multirow}
\usepackage{latexsym}
\usepackage{fullpage}
\usepackage{times}
\usepackage[standard]{ntheorem}
\usepackage{algorithmic}
\usepackage{algorithm}
\usepackage{capt-of}
\usepackage{listings}
\usepackage{url}
\usepackage{setspace}

\usepackage[    %pdftex,
        bookmarks=true,
	colorlinks=true,
        %bookmarksnumbered=false,
        bookmarksopen=true,
        bookmarksopenlevel=2,
        linkcolor=black,
        filecolor=black,
        anchorcolor=black,
        urlcolor=black,
        citecolor=black,
        pdfauthor={},
        pdfkeywords={},
	breaklinks=true,
        linkbordercolor=white,
        pdftitle={Titel der Arbeit}
        ]{hyperref}

\usepackage{vmargin}
\setmarginsrb
{25mm}{25mm} %% left top
{20mm}{25mm} % right bottom
{14pt}{20pt} % header: hight, gap
{0pt}{10mm}  % foot: hight, top

\parindent=0.4cm

\author{Student name}
\title{Master Thesis\\
\vspace{1cm}Title of the Thesis}

\date{
\begin{tabular}{r l}
Date:       & date \\
Supervisor: & name\\
\end{tabular}}

\publishers {Ruhr-University Bochum, Germany\\\vspace{0.5cm}
\begin{center}
  \includegraphics [height=5.5cm]{images/Label_RUB_WEISS-BLAU_srgb.jpg}
\end{center}
Chair for Network and Data Security\\
Prof. Dr. J\"org Schwenk\\
Homepage: \url{www.nds.rub.de}
\vspace{-4cm}}


%% Define fancy headers
\pagestyle{fancyplain}
\renewcommand{\chaptermark}[1]{\markboth{#1}{#1}} % remember chapter title
\renewcommand{\sectionmark}[1]{\markright{\thesection\ #1}} % section number and title
\lhead[\fancyplain{}{\thepage}]{\fancyplain{}{\sl \rightmark}}
\rhead[\fancyplain{}{\sl \leftmark}]{\fancyplain{}{\thepage}}
\cfoot{}
\renewcommand{\headrulewidth}{0pt}

%% here your document starts
\begin{document}

%% title page --- made out of expressions defined above
\maketitle

%% switch to roman paginating for the acknowledgements, table of contents etc.
\pagenumbering{roman}

%% include acknowledgement ("erklaerung.tex" is the one from RUB)
%
%
% Erklaerung
%
%

\section*{\begin{center}\LARGE \bf Erkl\"arung \end{center}}
Ich erkläre, dass das Thema dieser Arbeit nicht identisch ist mit dem Thema einer von mir
bereits für ein anderes Examen eingereichten Arbeit.
Ich erkläre weiterhin, dass ich die Arbeit nicht bereits an einer anderen Hochschule zur
Erlangung eines akademischen Grades eingereicht habe.\\

Ich versichere, dass ich die Arbeit selbständig verfasst und keine anderen als die angegebenen
Quellen benutzt habe. Die Stellen der Arbeit, die anderen Werken dem Wortlaut oder dem
Sinn nach entnommen sind, habe ich unter Angabe der Quellen der Entlehnung kenntlich
gemacht. Dies gilt sinngemäß auch für gelieferte Zeichnungen, Skizzen und bildliche
Darstellungen und dergleichen.\vspace{32pt}\\



\begin{table}[htbp]
\begin{center}
\begin{tabular}{lll}
\rule{0.4\textwidth}{0.4pt} & \hspace{42pt} &\rule{0.4\textwidth}{0.4pt} \\ 
Ort, Datum & \hspace{42pt} & Unterschrift
\end{tabular}
\end{center}
\end{table}


\newpage
%\begin{abstract}
 \begin{center}\textbf{\huge{Acknowledgements}}\end{center} \vspace{1cm}
The author would like to thank...
\clearpage
%\end{abstract}

\pagebreak


%% optional: include abstract
\begin{abstract}
\begin{center}\textbf{\huge{Abstract}}\end{center}

The abstract goes here.

\vspace{10mm}
KEYWORDS: some key words
\end{abstract}
\clearpage

\pagebreak

%% generate table of contents
\setcounter{tocdepth}{1}

\tableofcontents

\clearpage
\addcontentsline{toc}{section}{List of Figures}
\listoffigures

\clearpage
\addcontentsline{toc}{section}{List of Tables}
\listoftables

%%
%% include all your chapters as .tex files,
%% each file contains sections \section{name of section},
%% subsections \subsections{...} and so on...
%%
\setcounter{secnumdepth}{3}

\thispagestyle{empty}

\onehalfspacing
\chapter{Caption of Chapter 1}
  \label{chapter:1}
\pagenumbering{arabic} %switches to arabic numbers for the rest of the text
  Description of this chapter.

\section{Position of figures}
\label{section:1.1}
You can include the image as usual. You can include the image as usual.
You can include the image as usual. You can include the image as usual.
You can include the image as usual. You can include the image as usual.

\begin{figure}[ht]
\centering{\includegraphics [width=80mm]{images/NDS_Logo_HKS_42.pdf}}
\caption{SubBytes(normal figure)} \label{fig:1}
\end{figure}

You can include the image as usual. You can include the image as usual.
You can include the image as usual. You can include the image as usual.
You can include the image as usual. You can include the image as usual.

\begin{wrapfigure}{r}{80mm}
\centering{\includegraphics [width=80mm]{images/NDS_Logo_HKS_42.pdf}}
\caption{SubBytes(Normal figure (align right))} \label{fig:2}
\end{wrapfigure}

You can wrap the figure (align right). You can wrap the figure (align 
right). You can wrap the figure (align right). You can wrap the figure
(align right). You can wrap the figure (align right).

You can wrap the figure (align right). You can wrap the figure (align 
right). You can wrap the figure (align right). You can wrap the figure
(align right). You can wrap the figure (align right).

You can wrap the figure (align right). You can wrap the figure (align 
right). You can wrap the figure (align right). You can wrap the figure
(align right). You can wrap the figure (align right).

\begin{wrapfigure}{l}{80mm}
\centering{\includegraphics [width=80mm]{images/NDS_Logo_HKS_42.pdf}}
\caption{SubBytes(Normal figure (align left))} \label{fig:3}
\end{wrapfigure}

You can wrap the figure (align left). You can wrap the figure (align 
right). You can wrap the figure (align left). You can wrap the figure
(align left). You can wrap the figure (align left).

You can wrap the figure (align left). You can wrap the figure (align 
right). You can wrap the figure (align left). You can wrap the figure
(align left). You can wrap the figure (align left).

You can wrap the figure (align left). You can wrap the figure (align 
right). You can wrap the figure (align left). You can wrap the figure
(align left). You can wrap the figure (align left).

\section{Multiple Figures and Tables Parallel}
\label{section:1.2}

Two tables parallel.

\begin{figure}[htb]
  \centering
  \begin{minipage}[b]{0.45\textwidth}\center
  \begin{tabular}{|c|c|c|}\hline
   A & B & C\\\hline
  \end{tabular}
  \captionof{table}{Caption of table}
  \label{tbl:1}
  \end{minipage}\hspace{0.05\textwidth}
  \begin{minipage}[b]{0.45\textwidth}\center
  \begin{tabular}{|c|c|c|}\hline
  D & E & F\\\hline
  \end{tabular}
  \captionof{table}{Caption of table}
  \label{tbl:2}
  \end{minipage}
\end{figure}

Two figures parallel.

\begin{figure}[htb]
  \centering
  \begin{minipage}[b]{0.45\textwidth}\center
  \centering{\includegraphics [width=1.0\textwidth]{images/NDS_Logo_HKS_42.pdf}}
  \caption{Caption of figure}
  \label{fig:4}
  \end{minipage}\hspace{0.05\textwidth}
  \begin{minipage}[b]{0.45\textwidth}\center
  \centering{\includegraphics [width=1.0\textwidth]{images/NDS_Logo_HKS_42.pdf}}
  \caption{Caption of figure}
  \label{fig:5}
  \end{minipage}
\end{figure}

One figure and one table parallel

\begin{figure}[htb]
  \centering
  \begin{minipage}[b]{0.45\textwidth}\center
  \begin{tabular}{|c|c|c|}\hline
  D & E & F\\\hline
  \end{tabular}
  \captionof{table}{Caption of table}
  \label{tbl:3}
  \end{minipage}\hspace{0.05\textwidth}
  \begin{minipage}[b]{0.45\textwidth}\center
  {\includegraphics [width=1.0\textwidth]{images/NDS_Logo_HKS_42.pdf}}
  \caption{Caption of figure}
  \label{fig:6}
  \end{minipage}
\end{figure}

\newpage
\section{Algorithms}
\label{section:1.3}

Algorithm

\begin{algorithm}
\caption{Key pair generation of ECDSA}
\label{algo-crypto-6}

\begin{algorithmic}
\REQUIRE valid $PARAM$
\ENSURE public key $Q$ and secret key $d$
\WHILE{\TRUE}
\STATE $d \Leftarrow_R [1, $n$-1]$
\STATE $Q = (x_{_Q}, y_{_Q}) \Leftarrow dG$.
\IF{$Q \neq O$} 
\RETURN $(Q, d)$
\ENDIF
\ENDWHILE
\end{algorithmic}
\end{algorithm}

\section{Citations}
\label{section:1.4}

artice \cite{artice}\\
manual \cite[Section Nr.]{manual}\\
book \cite[Section Nr.]{inbook}\\
misc \cite[Section Nr.]{misc}\\
phdthesis \cite[Section Nr.]{phdthesis}\\
inproceedings \cite[Section Nr.]{inproceedings}\\
standard \cite[Section Nr.]{standard}\\
techreport \cite[Section Nr.]{techreport}\\









\chapter{Conclusions} 
  \label{chapter:con}
  Conclusions.

\singlespacing
%% include appendix
\appendix
\chapter{Caption of Appendix}
\label{appendix}
\section{Title of the first appendix section}
\label{sec:a-1}
You can embed the java source code as follows:

\lstset{language=java}
\begin{lstlisting}
public boolean verify(byte[] signature){
 return true;
}
\end{lstlisting}

\section{Caption of the second appendix section}
\label{sec:a-2}
BlaBla Blabla ...




%% generate bibliography with bibtex, the bibfile here is "biblist.bib"
%% use alphanumerical style
\clearpage
\addcontentsline{toc}{chapter}{Bibliography}
%
% ---- Bibliography ----
%
\bibliographystyle{IEEEtran}
%
%\IEEEtriggeratref{17}
\bibliography{IEEEabrv,ownbibfile}

%% this is the end ;)
\end{document}
